\chapter{Troubleshooting by symptom}\label{ch:troubleshooting}

Start from what the machine printed. The text below is the text ATE prints, copied from the code,
so searching for it finds this chapter.

\section{It will not start}

\textbf{The config is not there.} ATE names the path it expected and offers to create a default
config or abort:

\begin{lstlisting}[language=ATEShell]
No config file at /etc/ate/config.json
\end{lstlisting}

If it aborts instead --- for example because nobody was there to answer, or you chose Abort --- run
\texttt{ate setup} (as root on Linux) or point ATE at a config with \texttt{--config}.

\textbf{The config is there but wrong.} Every refusal names the file and the key:

\begin{lstlisting}[language=ATEShell]
unknown config key 'defnitions' in /etc/ate/config.json
config file /etc/ate/config.json: 'log.level' must be a string, got 3
config file /etc/ate/config.json: 'format' is required (expected 1)
config file /etc/ate/config.json: unknown 'format' value 2 (expected 1)
config file /etc/ate/config.json: 'database.type' must be 'sqlite', got 'postgres'
config file /etc/ate/config.json: 'ui.layout' must be 'staged' or 'threePane', got 'wide'
config file /etc/ate/config.json: 'devices.receiveBufferBytes' must be a positive integer, got 0
config file /etc/ate/config.json: 'reports.directory' has unknown placeholder '<datat>'
invalid JSON in config file /etc/ate/config.json: Expecting ',' delimiter: line 12
\end{lstlisting}

A key you did not write is not an error: missing keys are filled with their defaults. Only a key
that is present and wrong, or a key that does not exist at all, is refused.

\textbf{The config cannot be written.} A config with no \texttt{station.uuid} makes ATE generate one
and try to write it back:

\begin{lstlisting}[language=ATEShell]
cannot write station.uuid to /etc/ate/config.json; add '3f6c...' by hand and restart
\end{lstlisting}

That is a permissions problem, not a bug: whoever launched ATE cannot write \texttt{config.json}. An
admin adds the UUID, or runs ATE with the rights to write it.

\textbf{The database is not the right version.} ATE stops and says what to run:

\begin{lstlisting}[language=ATEShell]
run `ate db migrate` to create or upgrade the database at /var/lib/ate/ate.db
this ATE is too old for the database at /var/lib/ate/ate.db; update ATE
\end{lstlisting}

The first is fixed by migrating. The second means the database was built by a newer ATE than the one
installed; there is no downgrade, so update ATE or restore a backup.

\textbf{A directory is not writable.} ATE creates the log and reports directories at launch and
refuses plainly if it cannot:

\begin{lstlisting}[language=ATEShell]
cannot create /var/lib/ate/logs: [Errno 13] Permission denied. Run 'ate setup' as admin, or point
the config (log.directory, reports.directory) at a writable location. Config: /etc/ate/config.json
\end{lstlisting}

\section{\texttt{ate setup} did not do what you expected}

\begin{lstlisting}[language=ATEShell]
ate setup: running unprivileged; skipping group/ownership/ACL setup (packagers run this elevated)
\end{lstlisting}

This is expected when \texttt{ate setup} is not elevated: the directories, the config and the
database are still created, and the ownership treatment is skipped. Finish with
\texttt{sudo ate setup}. A real failure looks like \texttt{cannot create \dots}, \texttt{cannot write
config to \dots}, \texttt{database migration failed at \dots} or \texttt{cannot apply permissions to
\dots}, each with an exit code of 1.

If ordinary users cannot write reports or the database after setup, they are not in the \texttt{ate}
group on Linux, or the \texttt{icacls} grant did not apply on Windows. Setup makes the data tree
group-writable and the config admin-only; it does not add anyone to the group.

\section{A Test Run will not start, or warns}

\textbf{Files do not match their checksums.} The Run screen lists each as \texttt{expected \dots,
found \dots} before Start, and the run records:

\begin{lstlisting}[language=ATEShell]
checksum mismatch for lib/scpi.lua (expected 3f2a..., got 9b1c...)
\end{lstlisting}

The run still runs. Either the release was edited after it was shipped, or the checksums were never
recorded: run \texttt{ate definition update} on the \texttt{.atedef}.

\textbf{A schedule is missing from the picker.} Check the scan warnings. A directory that does not
exist warns; a broken definition is listed with its error and the rest of the scan is unaffected. A
schedule that raises while loading is broken on its own, not the whole release.

\textbf{A Test Report is not where you expected.} Reports are written under \texttt{reports.directory}
in a subdirectory named after the UUT Type, with the UUT Type, serial, local start time and result in
the file name. A name already taken is never overwritten: a second report for the same run gets
\texttt{-2}, then \texttt{-3}, and a name is considered taken if either the HTML or the PDF of that
name exists.

\section{There is no PDF}

The HTML Test Report is always written. If there is no PDF beside it, look in the application log.

\begin{lstlisting}[language=ATEShell]
Test Report PDF skipped for 0182...: ATE is privileged and the choice was No, so nothing was
rendered without Chromium's sandbox. HTML report: /var/lib/ate/reports/RailGun/..._PASS.html
\end{lstlisting}

ATE was run as root or an administrator, and the answer to "render without Chromium's sandbox?" was
no (or nobody answered). Run ATE as an ordinary user, or run it with
\texttt{--allow-unsandboxed-pdf} if giving up that isolation on this machine is acceptable to you.

A PDF that opens but is empty is a font problem, not a rendering problem: Debian and its containers
ship no fonts, and with none Chromium writes a valid PDF with no text in it. Install
\texttt{fontconfig} and \texttt{fonts-dejavu-core}. \texttt{ate \_selftest --run} says this in so
many words when it fails:

\begin{lstlisting}[language=ATEShell]
the PDF has no fonts: 1742 bytes with no text in it. Install fontconfig and a font (for example
fonts-dejavu-core) so QtWebEngine has something to draw with
\end{lstlisting}

If the helper itself fails you get one of these in the log, and the HTML report is still there:

\begin{lstlisting}[language=ATEShell]
PDF helper exited with code 1: ate _pdf: refusing to render as root or an administrator without --unsandboxed
PDF helper crashed or timed out after 60 s
PDF helper wrote no PDF
PDF helper exited with code 1: pdf helper error: paper must be 'A4' or 'Letter', got 'Legal'
\end{lstlisting}

The first appears only if you run \texttt{ate \_pdf} yourself: the parent always tells the helper
whether the sandbox was declined or agreed. The last is a \texttt{reports.paperSize} typo --- the
config does not check that key.

\section{A window stopped updating}

A status bar line appears, and stays:

\begin{lstlisting}[language=ATEShell]
Internal error: this window may have stopped updating. A Test Run in progress keeps recording. Details: /var/lib/ate/logs/ate-app-2026-11-01T09-02-31.log
\end{lstlisting}

The traceback is in that file, under \texttt{Unhandled exception on thread}. A Test Run that was in
progress keeps recording: the record is written by the worker, not the window, so trust the database
even when the screen is frozen. The details line is missing when the write itself failed, which says
the log directory is the thing to fix.

\section{A run is slow, or an instrument is still busy}

Two facts about Devices that an administrator needs, both established by the suite rather than
guessed.

\textbf{ATE takes no exclusive lock on a serial port on Linux.} Another program can open the same
\texttt{tty} while a Test Run is using it. The suite tried to prove the opposite with
\texttt{serial.Serial(exclusive=True)} and found the open succeeded even while the port was held, so
sharing is the behaviour, and two schedules or two programs on one port are your problem to avoid.

\textbf{An abort frees the worker, not the instrument.} When the Operator aborts a run, the worker
stops waiting for a Device call in about 50 ms and moves on; the call itself keeps running on a
daemon thread and holds its transport until it finishes. What can be left behind:

\begin{itemize}
  \item a connect or a DNS lookup (TCP, UDP, a \texttt{socket://} serial URL, VISA), bounded by the
    timeout the schedule gave or the operating system's own limit;
  \item \texttt{viRead} or \texttt{write\_raw} on anything other than pyvisa-py's SOCKET session ---
    NI-VISA, GPIB, USB, VXI-11 --- for as long as the call's timeout says;
  \item a pyserial write: the timeout does not say how much got out;
  \item \texttt{ATE.close} on a stuck port, which the operating system can hold for a long time
    (Linux can hold a \texttt{tty} close while output drains, up to 30 seconds by default).
  \item At the end of an aborted run all such closes together are given at most three seconds, with
    one WARNING naming each Device still closing.
\end{itemize}

A Post-test Action on the same Device waits for a call left behind (a second abort ends that wait),
so a run's end can run past ten seconds. A new run that opens the same instrument while an abandoned
call is still running may find it busy. Nothing about this is a fault to fix: it is what an abort
means, and it is why the warning line \texttt{was still closing when the aborted run ended} is not a
bug.

Three things cannot be interrupted at all, whatever the Operator presses: a SQLite commit already in
progress, one Lua C library call such as a \texttt{string.rep} or a large pattern match, and the
VISA attribute calls on the worker.

One stall has been seen and not explained. A Test Run on the Windows CI runner missed \texttt{FINISHED}
and its abort took over ten seconds; the evidence points at store commits stalling on that runner's
disk --- every commit in WAL mode is a synced commit at SQLite's \texttt{synchronous=FULL} --- but
nothing measured the disk, so the cause is a belief, not a finding. If it ever happens again, the
test report now prints each live worker's thread stacks, what it has recorded and the store's write
times, which separates a slow disk from a fault nobody has found yet.

\section{Runs closed as \texttt{INTERRUPTED}}

At launch ATE closes any run with no end time --- a crash, a kill, a power cut --- as
\texttt{INTERRUPTED} with result \texttt{FAIL}, and writes one line to the application log:

\begin{lstlisting}[language=ATEShell]
Closed 2 interrupted run(s) at launch.
\end{lstlisting}

That is the record saying what it knows: the run did not finish. It is not recoverable, and it is
not a sign that ATE is broken.

\cando{take a message an administrator sees and name the cause and the fix, and say when an abort
leaves an instrument busy rather than releasing it.}
